\section{One two-moment obstruction for the quadratic boundary cases}
\label{sec:quadratic-center-trace}

\begin{theorem}\label{thm:quadratic-center-trace}
The rings $R7\_4\_106$, $R7\_4\_2155$ and $R7\_3\_146$
admit no complex fusion categorification.
\end{theorem}

The proof combines full forgetful-row Galois transport, equality in the
modular compression bound, and a finite hull for two root-of-unity
moments. Four fixed integer inequalities suffice for all cases.
Proposition~\ref{prop:three-case-galois-pu} permits us to work with a
pseudo-unitary Galois conjugate and its canonical spherical structure;
no unitarity assumption on the original category is made.

\subsection{Source dimensions and complete row transport}
Write $d_i=a_i+b_i\sqrt r$ and $D=\sum_i d_i^2=48\eta$.
In source order the dimension vectors are
{\small
\[
\begin{array}{c|l|c}
\text{ring}&(d_0,\ldots,d_6)&\eta\\\hline
106&(1,1,4+2\sqrt3,7+4\sqrt3,7+4\sqrt3,4+2\sqrt3,6+4\sqrt3)&7+4\sqrt3\\
2155&(1,4+2\sqrt3,4+2\sqrt3,1,6+4\sqrt3,7+4\sqrt3,7+4\sqrt3)&7+4\sqrt3\\
146&(1,1,4+2\sqrt6,5+2\sqrt6,5+2\sqrt6,4+2\sqrt6,6+2\sqrt6)&5+2\sqrt6
\end{array}
\]
}
Here $\eta\overline\eta=1$. The formal codegrees are
$D,\overline D,4,8,8,8,12$ in the first two cases and
$D,\overline D,4,6,6,8,12$ in the third. Since the rings are
commutative, the seven distinct constituents $A_E$ of $I(1)$ have
multiplicity one, twist one and dimensions $D/f_E$
\cite[Theorems~2.5 and~2.13]{OstrikRank3}.

For a center simple $A$, let $v(A)_i=[F(A):X_i]$. Induction gives
\begin{equation}\label{eq:quadratic-center-gram}
 B_{ij}=\sum_{k,l}N_{ki}^lN_{kj}^l
       =\sum_A v(A)_i v(A)_j.
\end{equation}
If $z=\dim A=\sum_i v_i d_i$, modular formal-codegree integrality
makes $D^2/z^2$, hence $D/z$, algebraically integral. Since
$D/48$ is a unit and $\mathcal O_{\mathbb Q(\sqrt r)}=\mathbb Z[\sqrt r]$,
every $z=a+b\sqrt r$ satisfies
\begin{equation}\label{eq:quadratic-center-divisibility}
             a^2-rb^2\mid48a,\qquad a^2-rb^2\mid48b.
\end{equation}

Normalize the center matrix by $S_{00}=1/D$. The source characters
and the induced-unit constituents satisfy
\[
 \chi_E(F(A))=\frac{S_{A,A_E}}{S_{0,A_E}},\qquad
 S_{0,A_E}=1/f_E
\]
by the definition preceding \cite[Theorem~2.13]{OstrikRank3}.
Modular Galois symmetry consequently gives
\begin{equation}\label{eq:quadratic-row-galois}
 \chi_E(F(\pi_\sigma A))
 =\epsilon_\sigma(A)\frac{f_E}{f_{\sigma E}}
                          \chi_{\sigma E}(F(A)).
\end{equation}
The integral vector of multiplicities in $I(1)$ satisfies $Sm=m$
\cite[Proposition~2.8]{OstrikRank3}. Its signed Galois image is
therefore itself: this support is preserved and its signs are positive.

We choose $\sigma$ to flip the dimension field and fix all other
source characters. This requires more than rationality of their
codegrees. Set $M=\sum_i w_iN_i$, with
\[
                  w=(1,3,11,37,101,307,997).
\]
Its characteristic polynomials are
\begin{align*}
106:\quad &(x-1860)(x+66)(x+494)\\[-2pt]
 &\hspace{8mm}\cdot(x^2-16448x-12738752)(x^2-124x-171388),\\
2155:\quad &(x-1314)(x-654)(x+1064)\\[-2pt]
 &\hspace{8mm}\cdot(x^2-19656x+889872)(x^2+1452x+527204),\\
146:\quad &(x-654)(x-184)(x+66)(x+234)(x+1852)\\[-2pt]
 &\hspace{8mm}\cdot(x^2-15896x+12501688).
\end{align*}
All roots are distinct. The simultaneous eigenspaces are one-dimensional
and defined over the corresponding eigenvalue fields. Apart from the
dimension field, the character fields are respectively
$\mathbb Q(\sqrt2)$, $\mathbb Q(\sqrt{-2})$ and $\mathbb Q$.
Their disjointness from $\mathbb Q(\sqrt r)$ proves that the required
cyclotomic automorphism exists.

Put $a(v)=\sum_i a_iv_i$, $b(v)=\sum_i b_iv_i$ and
$\mathbf b=(b_i)$. Character orthogonality makes $\mathbf b$ vanish
on every character except the two dimension characters. Evaluation
on that pair in \eqref{eq:quadratic-row-galois} gives the integral
involution
\begin{equation}\label{eq:quadratic-row-transport}
 T(v)=\begin{cases}
 v+\mathbf b(4a(v)-7b(v))/4,&r=3,\\
 v+\mathbf b(2a(v)-5b(v))/2,&r=6.
 \end{cases}
\end{equation}
Every actual row has partner $T(v)$ or $-T(v)$, with a common sign
making all coordinates nonnegative. Mixed-sign images are impossible.
The multiplicities of paired row types agree; we do not assume that
the Galois permutation of individual center simples is an involution.

In particular, besides the unit row, $I(1)$ contains the forced row
\[
 (1,0,2,4,4,2,4),\quad(1,2,2,0,4,4,4),\quad
 (1,0,2,2,2,2,2),
\]
respectively, of dimension $\eta^2$. The five other nonunit rows have
their known dimensions, sum $B_{0,*}-e_0$ after including the forced
row, and positive semidefinite residual Gram matrix. Exhaustive bounded
integer enumeration gives $194,194,1323$ blocks. Every bound comes
from \eqref{eq:quadratic-center-gram}; no putative center is supplied
as input.

\subsection{Equality restricts the nonrational twist coordinates}
Let $q=12$ for $r=3$ and $q=24$ for $r=6$, and let $J$ be the
Galois image of the central unit. Every automorphism fixing the
dimension field fixes the central unit: all codegree ratios in
\eqref{eq:quadratic-row-galois} are then one, so its forgetful row
is $e_0$. Such a center simple lies in $I(1)$, where the dimension-one
constituent is unique. Those flipping the field send the unit to $J$.
Thus $w=e_{\mathbf1}-e_J$ has diagonal-torus character the primitive
quadratic character $\chi_q$. On the known twist-one coordinates,
\[
 \operatorname{proj}_{I(1)}Sw
   =\frac{\eta^{-1}-\eta}{48}w=-\frac1{\sqrt q}w.
\]
This attains the full compression bound. The equality argument in
Section~\ref{sec:three-case-galois-pu} therefore places the cyclic
representation generated by $w$ at level $q$.

We also need the sign of the primitive twist modes. More generally,
if $PSw=\delta w/\sqrt q$ with $\delta=\pm1$, a nonzero primitive
$q$-th twist mode has exponent $k$ satisfying $\chi_q(k)=\delta$.
Here is the local Fourier refinement. For a prime-power conductor
$q_i$, the equality subspace is spanned by
\[
 F_i(p,u)=\chi_i(u)[u\text{ a unit}],\qquad
 G_i(p,u)=\chi_i(-p)[u=0].
\]
The compressed Weyl operator sends $F_i$ to $G_i$ and $G_i$ to
$\chi_i(-1)F_i/q_i$. If $\tau_i$ is the local quadratic Gauss sum,
an eigenvector is $F_i+c_iG_i$ with
$c_i=\varepsilon_i\chi_i(-1)\tau_i$ and eigenvalue
$\varepsilon_i\tau_i/q_i$. On a unit-$u$ orbit, put $x=p/u$.
The primitive Fourier coefficient of its Weyl image is proportional to
\[
 \chi_i(-1)\tau_i\bigl(\chi_i(k)+\varepsilon_i\bigr).
\]
It vanishes unless $\chi_i(k)=\varepsilon_i$.
The Chinese-remainder identity
$\sqrt q=(\prod_i\chi_i(q/q_i))\prod_i\tau_i$ shows that the local
signs combine to $\chi_q(k)=\delta$. The assertion is preserved
under linear combinations in the common global eigenspace.

The coordinate of $Sw$ at a center simple of dimension $z$ is
$z/D-\sigma(z/D)$. Consequently
\begin{equation}\label{eq:quadratic-saturated-twist}
 z/D\notin\mathbb Q\quad\Longrightarrow\quad
 \theta^q=1,\qquad
 \chi_q(k)=-1\text{ if }\theta=e^{2\pi ik/q}\text{ is primitive}.
\end{equation}
Moreover, modular Galois symmetry sends twists to their $u^2$-th
powers. Every unit modulo $12$ or $24$ squares to one. The paired
nonrational rows therefore have equal twists, so their multiplicities
agree separately for each twist. Rational ratios are fixed by $T$ and
require no such restriction on the order of their twists.

\subsection{A finite two-moment hull with no order bound}
Let $t(\theta)$ be normalized relative trace to $K=\mathbb Q(\sqrt r)$,
taken from any cyclotomic extension containing $K$ and $\theta$.

\begin{lemma}\label{lem:quadratic-two-moment-hull}
For every root of unity, the pair $(t(\theta),t(\theta^2))$ belongs
to the convex hull of the following nine coefficient vectors, where
$(a,b,c,d)$ means $(a+b\sqrt r,c+d\sqrt r)$:
\begin{gather*}
 (1,0,1,0),\quad(-1,0,1,0),\quad
 (1/2,0,-1/2,0),\quad(-1/2,0,-1/2,0),\quad(0,0,-1,0),\\
 (0,h,s,0),\quad(0,-h,s,0),\quad
 (0,0,0,h),\quad(0,0,0,-h),
\end{gather*}
with $(h,s)=(1/2,1/2)$ for $r=3$ and $(1/4,0)$ for $r=6$.
Under \eqref{eq:quadratic-saturated-twist}, the first five vertices
and $(0,-h,s,0)$ suffice.
\end{lemma}
\begin{proof}
If the root order $n$ is not divisible by the quadratic conductor $q$,
the trace is the Ramanujan value $\mu(n)/\varphi(n)$. For $q\mid n$,
the rational contribution vanishes. The induced quadratic Gauss sum
is nonzero precisely when $n=qm$, where $m$ is squarefree and coprime
to $6$. Put $c=h\sqrt r$. The corresponding pairs are
\[
 \left(\frac{\pm c\mu(m)\chi_q(m)}{\varphi(m)},
                    \frac{\mu(m)}{2\varphi(m)}\right)
 \quad(q=12),\qquad
 \left(\frac{\pm c\mu(m)\chi_q(m)}{\varphi(m)},0\right)
 \quad(q=24).
\]
For $m=1$ these are the primitive-conductor vertices. Otherwise
$w=1/\varphi(m)\leq1/4$. Give the appropriate primitive vertex weight
$w$. When $q=12$ and $\mu(m)=-1$, the remaining weight supplies
the vertical correction $-w/(1-w)\in[-1,1]$; otherwise it supplies
zero. The whole vertical segment $(0,[-1,1])$ is in the rational
five-vertex hull.

A nonzero irrational second moment with zero first moment requires
$n=2qm$ with the same restrictions on $m$ and lies on the segment
between the two primitive-$2q$ vertices. In the rational cases, odd
squarefree order gives $(u,u)$, with $u\in[-1/2,1]$, and twice odd
squarefree order gives $(-u,u)$. Orders divisible by four give
$(0,u)$, $u\in[-1,1]$. The other cases give zero. These are all
in the displayed hull. Finally, for order dividing $q$, the primitive
negative-character vertex and the rational five-vertex hull cover
all possibilities in \eqref{eq:quadratic-saturated-twist}.
\end{proof}

\subsection{Pivotal signs and four exact inequalities}
For arbitrary $X$ and self-dual $Y$, the pivotal form
$j_X\otimes f_Y\otimes\mathrm{id}_{X^*}$ on $XYX^*$ has sign
$\nu_2(Y)$. Every self-dual constituent of odd multiplicity therefore
has that same sign: a nondegenerate alternating form cannot exist
on its odd-dimensional multiplicity space. The same palindromic
argument applies to $X^3$ for self-dual $X$. Applied to the source
products, these relations leave exactly
\[
\begin{array}{c|c}
106&(1,1,\epsilon,1,1,\epsilon,1),\quad\epsilon=\pm1\\
2155&(1,0,0,1,1,1,1)\\
146&(1,1,1,1,1,1,\epsilon),\quad\epsilon=\pm1.
\end{array}
\]
The zeros are the nonselfdual indicators. In the second ring, label
$4$ occurs once in $X_1X_1^*=X_1X_2$, which fixes its indicator
to $+1$.

For every nonunit $X_i$, the induction trace identities are
\[
 \operatorname{Tr}(\theta\mid I(X_i))=0,\qquad
 \operatorname{Tr}(\theta^2\mid I(X_i))=D\nu_2(X_i)
\]
\cite[Theorems~2.5 and~2.7]{OstrikRank3}. Let $\ell$ run over the
six nonunit constituents of $I(1)$, write $z_\ell=\ell\cdot d$,
and put $k_i=\sum_\ell\ell_i z_\ell$. After their subtraction,
the residual $6\times6$ Gram matrix is
$R_{ij}=B_{ij}-\sum_\ell\ell_i\ell_j$, $1\leq i,j\leq6$.
Taking relative traces gives a system with $21$ Gram equations,
$12$ first-trace equations and $12$ second-trace equations.

For a paired orbit $\mathcal O$ of residual rows $(v,z)$ and a hull
vertex $(a,b,c,d)$, its column is
\begin{equation}\label{eq:quadratic-moment-column}
 \begin{pmatrix}
 4\sum_{v\in\mathcal O}\operatorname{vech}(vv^{\mathsf T})\\
 4\sum_{(v,z)\in\mathcal O}(v_i z(a+b\sqrt r))_{i=1}^6\\
 4\sum_{(v,z)\in\mathcal O}(v_i z(c+d\sqrt r))_{i=1}^6
 \end{pmatrix},
 \qquad
 b=4\begin{pmatrix}
 \operatorname{vech}(R)\\(-k_i)_{i=1}^6\\
 (D\nu_2(X_i)-k_i)_{i=1}^6
 \end{pmatrix}.
\end{equation}
The last two blocks are expanded in the coefficient basis $1,\sqrt r$.
The lower-triangular order is $(i,j)$ with $i$ increasing and then
$0\leq j\leq i$, using zero-based indices for the six residual
coordinates. All columns and targets are integral. A fixed row is
counted once; a two-row orbit is counted twice through the displayed
sum. Actual center objects give a nonnegative real combination of
these columns, using the same hull vertex on Galois-paired rows.

Every admissible $v$ is bounded by the Gram diagonal and satisfies
\eqref{eq:quadratic-center-divisibility}, the entrywise Gram bounds,
and $vv^{\mathsf T}\leq R$. The latter test includes singular $R$:
if $P=R_{I,I}$ is an invertible principal block, it is equivalent to
\[
 R_{*,I}P^{-1}v_I=v,\qquad v_I^{\mathsf T}P^{-1}v_I\leq1.
\]
The verifier applies these tests in exact rational arithmetic.
Some blocks are excluded already because $\operatorname{vech}(R)$
does not lie in the finite-field span of the paired Gram columns.
For every remaining block and every allowed indicator vector,
the four covectors in Table~\ref{tab:quadratic-separators} satisfy
\begin{equation}\label{eq:quadratic-farkas}
               y\cdot c\geq0\quad\hbox{for every allowed column},
               \qquad y\cdot b<0.
\end{equation}
Thus no nonnegative combination can equal the target. The first
covector covers both signs for $106$; the second covers $2155$;
the final two cover the positive and negative signs for $146$.

\begin{center}
\begin{tabular}{lrrrrr}\toprule
Ring&Low blocks&Modular exclusions&Trace blocks&Trace systems&Covectors\\\midrule
$106$&194&21&173&346&1\\
$2155$&194&21&173&173&1\\
$146$&1323&235&1088&2176&2\\\bottomrule
\end{tabular}
\end{center}
The $1088$ blocks for $146$ include repeated targets; all are checked,
not just the $1072$ different right-hand sides for each sign.
The source-bound verifier \texttt{verify\_quadratic\_center\_trace.py}
reconstructs the census IDs, immutable tensor hashes, codegrees,
character fields, low blocks, residual rows and every column of
\eqref{eq:quadratic-moment-column}. It uses only standard-library
integer and rational arithmetic; neither numerical feasibility nor
saved candidate lists are trusted. Independent ordinary and optimized
replays verify the complete coverage and every inequality
\eqref{eq:quadratic-farkas}. This excludes every pseudo-unitary center
possibility and proves Theorem~\ref{thm:quadratic-center-trace}.

\begin{table}[p]
\centering\small
\caption{The four universal covectors in the coordinate order of \eqref{eq:quadratic-moment-column}. Rows 1--21 are Gram coefficients, 22--33 first-trace coefficients, and 34--45 second-trace coefficients.}\label{tab:quadratic-separators}
\begin{tabular}{r r r r r}\toprule
Coordinate & $106$ & $2155$ & $146,+$ & $146,-$\\\midrule
1 & 1004 & 3 & 1006 & 1006\\
2 & 512 & -18 & 1000 & 1000\\
3 & 382 & 3 & 420 & -494\\
4 & 1000 & 1000 & 1000 & 1000\\
5 & -262 & 1000 & -377 & -291\\
6 & 34 & 1012 & 53 & 228\\
7 & 1000 & -148 & 1000 & 1000\\
8 & -290 & -148 & -319 & -594\\
9 & 51 & 1000 & 19 & 182\\
10 & 42 & 496 & 83 & 118\\
11 & 1000 & -383 & 1000 & 1000\\
12 & 1000 & -383 & 828 & -1000\\
13 & -259 & 1000 & -377 & -291\\
14 & -293 & 781 & -319 & -594\\
15 & 626 & 97 & 420 & -494\\
16 & 1000 & -293 & 606 & -1000\\
17 & -286 & -293 & -1000 & 559\\
18 & 339 & 1000 & 494 & 135\\
19 & 375 & 658 & 406 & 1000\\
20 & -773 & 17 & -1000 & 559\\
21 & -246 & 90 & 699 & 1006\\
\midrule
22 & -572 & 23 & 399 & 452\\
23 & 1000 & -176 & -1000 & -1000\\
24 & -580 & 23 & 174 & 303\\
25 & 1000 & -176 & -477 & -1000\\
26 & 141 & -739 & -316 & 9\\
27 & -247 & 1000 & 795 & 0\\
28 & -140 & 729 & 135 & 421\\
29 & 247 & -1000 & -318 & -1000\\
30 & -580 & 578 & -417 & -497\\
31 & 1000 & -1000 & 1000 & 1000\\
32 & 579 & 574 & 423 & -236\\
33 & -1000 & -1000 & -1000 & 1000\\
\midrule
34 & -774 & -113 & -698 & -418\\
35 & 1000 & 74 & 1000 & 1000\\
36 & 436 & -142 & 454 & -569\\
37 & -763 & 125 & -1000 & 1000\\
38 & -152 & -1000 & 406 & -416\\
39 & 259 & 975 & -1000 & 1000\\
40 & -146 & 718 & 407 & -388\\
41 & 249 & -1000 & -1000 & 1000\\
42 & 434 & 576 & 454 & -569\\
43 & -759 & -1000 & -1000 & 1000\\
44 & 594 & 480 & 198 & -142\\
45 & -1000 & -832 & -606 & 1000\\
\bottomrule\end{tabular}
\end{table}

